% ===================================================================
% chapters/02c_literature_methods.tex -- writer L3
% Sections of the literature chapter (no \chapter): methods, mass and
% permittivity data, compact/TCAD modelling literature, DFT, secondary
% sources and the Silvaco ATLAS 5.28.1.R manual.
% ASCII only. Page numbers follow CONVENTIONS section 5.
% ===================================================================

\section{Stokey et al.\ (2021): dielectric constants and effective mass of cubic \InO}
\label{sec:lit-stokey2021}

\paragraph{Summary.}
\citet{Stokey2021} determine the infrared dielectric function of bulk single-crystal \InO{} by
ellipsometry combined with density-functional perturbation theory, and measure the conduction-band
effective mass by the optical Hall effect. They review earlier work, including the strongly
nonparabolic conduction band and the zero-density mass parameter $\mstar(0) = (0.18 \pm 0.02)\,m_e$
of Feneberg et al.\ \citep[p.~225102-2]{Stokey2021}, and the static dielectric constant of Hamberg and
Granqvist \citep[p.~225102-3]{Stokey2021}. From the phonon-mode parameters (Table~II) and the
Lyddane--Sachs--Teller relation they obtain $\varepsilon_\infty = 4.05 \pm 0.05$ and
$\varepsilon_{\mathrm{DC}} = 10.55 \pm 0.07$ \citep[p.~225102-12, Table~II]{Stokey2021}, and from the
optical Hall data an effective mass $\mstar = (0.208 \pm 0.006)\,m_e$ at a free-electron density of
$2.81\times10^{17}\cmcu$ \citep[pp.~225102-12--13]{Stokey2021}.

\begin{usedbox}
\textbf{Items used from \citet{Stokey2021}.}\par\smallskip
\footnotesize
\begin{tabularx}{\linewidth}{@{}>{\raggedright\arraybackslash}p{1.8cm}>{\raggedright\arraybackslash}p{1.7cm}>{\raggedright\arraybackslash}p{1.6cm}>{\raggedright\arraybackslash}X>{\raggedright\arraybackslash}p{1.35cm}>{\raggedright\arraybackslash}X@{}}
\toprule
Quantity & Value & Page & How it enters the model & Section & Caveat \\
\midrule
Bulk $\mstar$ & $0.208\,m_0$ & pp.~225102-12--13 & constant term of $\mstar(t) = 0.208 + 0.1311\,t^{-1.4121}$ ($m_0$); sets $\Nc(t) = 2.509\times10^{19}(\mstar/m_0)^{3/2}\cmcu$ = 3.27, 2.55, 2.44, $2.40\times10^{18}$ for 2.0, 6.3, 13.2, 31.8\,nm & \cref{sec:par-mstar,sec:par-nc} & pure crystalline \InO{} at low density; W and disorder may raise $\mstar$; tested by $\mstar \times1.3$ (\run{0050}, \run{0051}) \\
$\varepsilon_{\mathrm{DC}}$ & 10.55 & p.~225102-12, Table~II & upper sensitivity bound of the channel permittivity (\run{0026}) & \cref{sec:par-eps-iwo} & bulk crystal; the model value is the C--V 9.30 \\
\bottomrule
\end{tabularx}
\end{usedbox}

\paragraph{Considered but not used.}
The zero-density extrapolation $0.18\,m_0$ and the nonparabolicity of the conduction band
\citep[p.~225102-2]{Stokey2021} were not used, because the model employs a parabolic band (ATLAS
continuum DOS). The effect of nonparabolicity is absorbed in the exponent of the confinement law,
which is smaller than 2, and is discussed in \cref{sec:der-confinement}. The high-frequency constant
$\varepsilon_\infty = 4.05$ (p.~225102-12) has no role in a static device simulation. The bulk
mobilities of the crystals (p.~225102-13) are not transferable to nanometre-thick
amorphous/nanocrystalline \IWO.

\paragraph{Contradictions.}
Three values of the bulk mass are available: $0.208$ (optical Hall, $n = 2.8\times10^{17}\cmcu$),
$0.18$ (zero-density extrapolation, Feneberg) and 0.17 (PBE, \citealp[p.~21540]{Lin2022}). Their
spread of $\sim$20\% is smaller than the $\times1.3$ sensitivity test, which shifted $\Vthcc$ by
$-44$\,mV at \SI{2}{nm} (\run{0050}: 0.637 versus 0.680\,V in \run{0038}) and by $-29$\,mV at
\SI{13.2}{nm} (\run{0051}: 0.027 versus 0.055\,V in \run{0040}).

\section{Pandey et al.\ (2025): analytical short-channel model of BEOL oxide TFTs}
\label{sec:lit-pandey2025}

\paragraph{Summary.}
\citet{Pandey2025} derive a two-dimensional analytical potential model for oxide-semiconductor
thin-film transistors with a Green's-function approach, motivated by the high carrier densities
($10^{19}$--$10^{20}\cmcu$) of \InO-type channels \citep[p.~1]{Pandey2025}. They compare the Green's
function and scale-length approaches and approximate a Gaussian trap distribution by an integrable
function \citep[p.~2]{Pandey2025}, then build the potential from Green's identity
\citep[p.~3]{Pandey2025}. The model is checked against Sentaurus simulations calibrated to \InO{}
data, in which the channel trap concentration is $10^{20}\cmcu$ at \SI{0.4}{eV} above the conduction-band
edge, with a Gaussian trap distribution and an intrinsic density of $5\times10^{19}\cmcu$
\citep[p.~5]{Pandey2025}; the trap density in the source/drain regions is set to $4\times10^{20}\cmcu$, and it stays fixed there in the later channel-trap tests
\citep[pp.~5, 7]{Pandey2025}.

\begin{usedbox}
\textbf{Items used from \citet{Pandey2025}.}\par\smallskip
\footnotesize
No number from this paper enters the model. It is used as a method reference: (i) a calibrated TCAD
treatment of \InO{} places a large trap density at the charge-neutrality-like level
$\approx$\SI{0.4}{eV} above $\Ec$ (p.~5), which is the same physical picture as the CNL argument used
as a consistency check in \cref{sec:par-chi}; (ii) Gaussian trap distributions are standard in
oxide-TFT modelling (p.~2), as used for the deep acceptor Gaussian and the interface sheet
(\cref{sec:par-deep,sec:par-dit}).
\end{usedbox}

\paragraph{Considered but not used.}
The short-channel analysis itself was not used. The target devices have $L = \SI{20}{\micro\metre}$,
so that two-dimensional short-channel effects are negligible and the ATLAS model is in effect a
long-channel simulation (\cref{sec:inp-geometry}). The trap magnitudes $10^{20}$--$4\times10^{20}\cmcu$ (pp.~5, 7)
describe undoped \InO{} and are not transferred.

\paragraph{Contradictions.}
The \HfO{} relative permittivity used in their example is 15 \citep[p.~2]{Pandey2025}, lower than the
19.57 measured on the target process (SI Fig.~S1b, via the Astra audit); the model uses 19.57
(\cref{sec:par-stack}).

\section{Wang X.\ et al.\ (2025): physics-based compact model of IGZO FETs}
\label{sec:lit-wangx2025}

\paragraph{Summary.}
\citet{WangX2025} propose a compact model for IGZO-channel FETs in which the shallow donors are
described by a Gaussian density of states that acts as positive charge in depletion
\citep[pp.~2390--2391]{WangX2025}. The charge of the Gaussian distribution requires a Gauss--Fermi
integral; they take over a published error-function approximation of it and replace the error
function by a tanh function so that the model can be written in Verilog-A
\citep[pp.~2391--2392]{WangX2025}; the electrostatics apply Gauss's law across the gate oxide
\citep[p.~2393]{WangX2025}, and device-to-device variability is analysed with Gaussian-distributed
oxide thickness, channel thickness and donor density \citep[p.~2396]{WangX2025}.

\begin{usedbox}
\textbf{Items used from \citet{WangX2025}.}\par\smallskip
\footnotesize
No number is used. The paper supports the modelling choice recorded in
\file{config/iwo_material_model.yaml}: a donor Gaussian at a fixed absolute energy (ATLAS NGD/EGD)
is the natural V2 replacement for the empirical $\Ndeff(t)$ law, because confinement would then
ionise or neutralise the donors automatically as $\Ec$ moves (\cref{sec:par-nd}). In V1 the donor
Gaussian is set to zero (NGD $= 0$, \prov{ASSUMED}).
\end{usedbox}

\paragraph{Considered but not used.}
The closed-form Gauss--Fermi approximation (pp.~2391--2392) was not used, because ATLAS integrates
the discretised DOS numerically and no analytical approximation is needed. The IGZO parameter values
are not transferable to \IWO.

\paragraph{Contradictions.} No contradiction relevant to the model was found.

\section{Anusha and Dwivedi (2024): review of TCAD and compact modelling of IGZO TFTs}
\label{sec:lit-anusha2024}

\paragraph{Summary.}
\citet{Anusha2024} review the simulation and compact modelling of amorphous IGZO transistors,
including the commercial structure editors of Sentaurus and Silvaco \citep[p.~2]{Anusha2024}. They
summarise work on W-doped InZnO (WIZO) transistors in which the field-effect mobility declines as the
W content increases \citep[p.~4, Fig.~2a]{Anusha2024}, and quote the standard relation between the
subthreshold swing and the bulk and interface trap densities,
$\mathit{SS} \propto (\kB T/q)\,(N_{ss}t_{ch} + \Dit)$ (full form with $\Cox$ in the source),
used to separate $N_{ss}$ and $\Dit$ by fixing one of them \citep[p.~5]{Anusha2024}. Typical IGZO
swings of about \SI{120}{mV/dec} are mentioned in the outlook \citep[p.~14]{Anusha2024}.

\begin{usedbox}
\textbf{Items used from \citet{Anusha2024}.}\par\smallskip
\footnotesize
\begin{tabularx}{\linewidth}{@{}>{\raggedright\arraybackslash}p{2.2cm}>{\raggedright\arraybackslash}p{1.8cm}>{\raggedright\arraybackslash}p{1.0cm}>{\raggedright\arraybackslash}X>{\raggedright\arraybackslash}p{1.35cm}>{\raggedright\arraybackslash}X@{}}
\toprule
Quantity & Value & Page & How it enters the model & Section & Caveat \\
\midrule
SS--trap relation & $N_{ss}t_{ch} + \Dit$ form & p.~5 & theory used to interpret the fixed-current SS and the trap-DOS inversion; not a deck input & \cref{sec:der-ss} & valid only where the traps are in equilibrium and uniformly distributed; bulk and interface terms are degenerate from one curve \\
W doping lowers mobility (WIZO) & qualitative & p.~4 & context for the low $\muband$ of \IWO{} vs \InO & \cref{sec:par-mu} & different host (InZnO) \\
\bottomrule
\end{tabularx}
\end{usedbox}

\paragraph{Considered but not used.} The IGZO parameter sets of the review were not used, because
they do not describe \IWO.

\paragraph{Contradictions.} No contradiction was found.

\section{Giannozzi et al.\ (2009): Quantum ESPRESSO and the unexecuted DFT workflow}
\label{sec:lit-giannozzi2009}

\paragraph{Summary.}
\citet{Giannozzi2009} describe Quantum ESPRESSO, an open-source suite for electronic-structure
calculations \citep[p.~1]{Giannozzi2009} based on density-functional theory, plane-wave basis sets and
norm-conserving, ultrasoft or projector-augmented-wave pseudopotentials, distributed under the GNU
General Public License \citep[p.~2]{Giannozzi2009}.

\begin{usedbox}
\textbf{Items used from \citet{Giannozzi2009}.}\par\smallskip
\footnotesize
The paper serves as a software reference only. The package contains a prepared QE slab workflow
(\file{qe_workflow/}) for computing the vacuum-aligned band edges and slab masses of \IWO{} directly;
its results would replace the pure-\InO{} proxies of Lin2022 and Si2021 and fix the absolute electron
affinity. The workflow was not executed, because \texttt{pw.x} is absent on this host (Windows and
WSL) and its installation required an interactive password (\file{docs/INPUT_INVENTORY.md}). No DFT
number in this report was computed by the project; all DFT numbers are literature PBE values
(\prov{DFT\_PURE\_IN2O3\_PROXY}).
\end{usedbox}

\section{Secondary sources that are not bundled locally}
\label{sec:lit-secondary}

The works in \cref{tab:lit-secondary} were not available as text, and none of them was read for this
report. Each appears only through the bundled record that cites it (with the printed page of that
record) or through the Astra audit notes (\file{audit/docs/paper_revision_20260911/defect_papers.md} and
\file{ald_paper.md} in the Astra package), and is labelled accordingly. Values that exist only in the
Astra audit are marked ``via Astra'' and could not be verified locally.

\begin{table}[htbp]
\centering
\footnotesize
\caption{Secondary (not bundled) sources: what is attributed to them, the citing record, and their use.}
\label{tab:lit-secondary}
\begin{tabularx}{\linewidth}{@{}>{\raggedright\arraybackslash}p{1.9cm}>{\raggedright\arraybackslash}X>{\raggedright\arraybackslash}p{3.4cm}>{\raggedright\arraybackslash}p{3.2cm}@{}}
\toprule
Key & Attributed content & Citing record & Use in the model \\
\midrule
\citet{JanuarSI2026} & MIM/MISM C--V permittivities: \IWO{} 9.30$\pm$0.04, \InO{} 9.03, \HfO{} 19.57 (Fig.~S1b); grain sizes (Fig.~S4); annealing series (Figs.~S7, S8); 2\,nm transfer curves at 338/358\,K (Fig.~S12) & captions printed in the main text, \citet[p.~11]{Januar2026}; values recorded by the Astra audit & $\epsIWO = 9.30$ and $\varepsilon_{\mathrm{HfO_2}} = 19.57$ (\prov{MEASURED\_TARGET\_DEVICE}, \cref{sec:par-eps-iwo,sec:par-stack}); Fig.~S12 is the first test of \run{0044}--\run{0047} \\
\citet{Fan2021} & bulk $\Eg = 3.05$\,eV and $\chi = 4.30$\,eV of a sputtered a-\IWO{} (Sec.~3; $\chi$ a composition-mixing estimate); deep-acceptor range 2.5--3.7$\times10^{16}$ (p.~9); $\Nc$ used as a fit knob; printed inconsistencies noted by Astra & via Astra (\file{defect_papers.md}) & $\Eg$, $\chi$ bulk priors (\prov{LITERATURE\_IWO}, \cref{sec:par-eg,sec:par-chi}); $\NGA = 5\times10^{16}$ (\prov{ASSUMED}, \cref{sec:par-deep}) \\
\citet{Hamberg1986} & static permittivity $\approx$8.9 of evaporated (Sn-doped) \InO{} films & \citet[p.~501]{Si2021}; \citet[p.~4]{Wang2022}; \citet[p.~168]{Si2022}; \citet[p.~225102-3]{Stokey2021} & lower bound of the permittivity range only \\
Feneberg et al.\ (no key) & zero-density mass $\mstar(0) = 0.18\pm0.02\,m_e$, nonparabolicity & \citet[p.~225102-2]{Stokey2021} & considered; 0.208 used \\
\citet{King2008} & charge-neutrality level above $\Ec$, surface electron accumulation in \InO & cited for TNL $\approx$0.4\,eV in \citet[p.~504]{Si2021} (ref.~26, listed p.~506) & consistency check only (\cref{sec:par-chi}) \\
\citet{Robertson2011} & limits to doping in oxides; CNL alignment & \citet[p.~504]{Si2021} (ref.~24); \citet[p.~5]{Wang2022} & consistency check only \\
\citet{Hinuma2019} & band alignment at surfaces and interfaces & listed as ref.~27 in \citet[p.~506]{Si2021} & none numeric \\
\citet{Ando1982} & surface-roughness scattering and the mobility roll-off $\mu_0/(1+\theta_r V_{\mathrm{ov}})$ & \citet[p.~7]{Januar2026} (ref.~39, listed p.~11) & form of the roughness factor; roll-off not inserted (\cref{sec:par-rough,sec:par-mu}) \\
\citet{Yoo2024} & ALD \IWO{} channel & ref.~14 of \citet[p.~10]{Januar2026}; reviewed in the Astra \file{ald_paper.md} & none numeric (process context) \\
\citet{Lee2018} & gate-current effects in oxide TFTs (title only; content not read here) & no bundled record cites it; listed via the Astra package bibliography & none numeric; attribution to be confirmed \\
\bottomrule
\end{tabularx}
\end{table}

\section{The Silvaco ATLAS 5.28.1.R user manual: every statement used by the decks}
\label{sec:lit-atlas-manual}

The decks (\cref{sec:num-deck}, full listings in \cref{ch:app-decks}) use a small set of native
ATLAS statements. \Cref{tab:lit-atlas} gives, for each statement group, the printed manual page(s)
where its syntax and parameters are defined, and what the decks set with it. Page numbers are the
printed page numbers of \file{atlas_users1.pdf} (ATLAS 5.28.1.R); the PDF page is equal to the
printed page at p.~179, 4 higher at p.~513 and 6 higher throughout the statement chapter
(pp.~1117--1647) (\file{digest/atlas_manual_pages.md}).

\begin{table}[htbp]
\centering
\footnotesize
\caption{ATLAS statements used by the decks, with the printed manual pages and the model content.}
\label{tab:lit-atlas}
\begin{tabularx}{\linewidth}{@{}>{\raggedright\arraybackslash}p{3.0cm}>{\raggedright\arraybackslash}p{2.3cm}>{\raggedright\arraybackslash}X>{\raggedright\arraybackslash}p{2.0cm}@{}}
\toprule
Statement (parameters) & Manual pages & What the decks set & Section \\
\midrule
MESH, X.MESH, Y.MESH & \citep[pp.~1117--1118, 1373--1377]{AtlasManual} & structured 2-D mesh; fine vertical spacing in the film; spacing factor 1.0 (0.7 in the checks \run{0028}, \run{0036}) & \cref{sec:num-mesh} \\
REGION & \citep[pp.~1548--1560]{AtlasManual} & TiN gate, \HfO{} \SI{15}{nm}, \AlO{} \SI{2}{nm}, \IWO{} film, 0.25\,nm interface layer, back-surface layer, air & \cref{sec:par-stack,sec:inp-geometry} \\
ELECTRODE & \citep[pp.~1207--1211]{AtlasManual} & gate, source, drain (Pd, 2\,$\mu$m overlap) & \cref{sec:par-contacts} \\
DOPING UNIFORM & \citep[pp.~1183--1193]{AtlasManual} & $\Ndeff(t)$ as uniform n-type background & \cref{sec:par-nd} \\
MATERIAL (EG300, AFFINITY, PERMITTIVITY, NC300, NV300, MUN, MUP, TAUN0, EGALPHA, TMUN) & \citep[pp.~1307--1355]{AtlasManual} & $\Eg(t)$, $\chi(t)$, $\epsIWO$, $\Nc(t)$, $\Nv$, $\mu_0(t) = \muband(t)(1-\Dsr/t)^2$, hole mobility 0.01, SRH lifetime & \cref{sec:par-eg,sec:par-chi,sec:par-nc,sec:par-mu} \\
Constant mobility temperature dependence (TMUN) & \citep[p.~179]{AtlasManual}; \citep[p.~513]{AtlasManual} & $\mu \propto (T/300)^{-\mathrm{tmu}}$ with tmu $= 1.5$ in the elevated-temperature runs \run{0044}--\run{0047} & \cref{sec:par-temperature} \\
CONTACT (WORKFUNCTION, RESISTANCE) & \citep[pp.~1145--1156]{AtlasManual} & gate work function \SI{4.70}{eV}; Ohmic S/D (no work function); lumped series resistance for \SI{31.8}{nm} only & \cref{sec:par-gatewf,sec:par-contacts} \\
MODELS (FERMI, SRH, TEMPERATURE, PRINT) & \citep[pp.~1426, 1433, 1444--1448, 1453, 1462, 1477--1478]{AtlasManual} & Fermi--Dirac statistics, SRH, lattice temperature (300\,K; 338.15/358.15\,K in predictions) & \cref{sec:num-solver,sec:par-temperature} \\
DEFECTS (CONTINUOUS, NUMA, NUMD, NTA, WTA, NGA, EGA, WGA, NGD, EGD, SIGTAE) & \citep[pp.~1167--1174]{AtlasManual} & acceptor tail $\NTA(t)$, $\WTA = 0.040$\,eV; deep Gaussian; interface Gaussian in the 0.25\,nm layer; back-surface donor Gaussian (31.8\,nm hypothesis); NUMA/NUMD = 96/48, 192/96, 384/192 & \cref{sec:par-tail,sec:par-deep,sec:par-dit,sec:der-dos-discretization} \\
INTERFACE (QF) & \citep[pp.~1252--1264]{AtlasManual} & fixed charge $\Qf$ at the \AlO/\IWO{} interface & \cref{sec:par-qf} \\
METHOD (NEWTON, TRAP, MAXTRAPS, ITLIMIT, CLIMIT, IR.TOL, CR.TOLER, XANDRNORM, CARRIERS) & \citep[pp.~1378--1390, 1395]{AtlasManual} & Newton solver, electrons only, bias-step reduction on failure & \cref{sec:num-solver} \\
SOLVE (INIT, VSTEP, VFINAL, NAME) & \citep[pp.~1587--1611]{AtlasManual} & stepped gate sweeps $-3 \rightarrow +3$\,V (steps 0.2 / 0.05 / 0.1\,V) at $\Vd = 0.7$\,V; ID--VD sweeps & \cref{sec:num-sweep} \\
LOG, SAVE, PROBE, OUTPUT, EXTRACT & \citep[pp.~1216, 1299--1305, 1504--1515, 1524--1540, 1566--1579]{AtlasManual} & terminal currents, structure files, band diagrams (EXTRACT is only listed in the ATLAS manual, p.~1216, and documented in the DeckBuild manual) & \cref{sec:num-deck} \\
\bottomrule
\end{tabularx}
\end{table}

Two facts about the simulator are relevant to the physics; both were established by runs rather than
taken from the manual. First, ATLAS 5.28.1.R ignores MOBILITY updates between SOLVE statements (V0
Claude run~0008, recorded in \file{config/iwo_material_model.yaml}); for this reason the
gate-dependent roll-off of \citet[p.~7]{Januar2026} could not be inserted. Second, the
constant-mobility temperature exponent tmu $= 1.5$ is active in the elevated-temperature runs
(SYNTHESIS correction), so that the registered temperature predictions exist in a phonon-like and an
MTR-like variant (\cref{sec:pred-temperature}).
